\chapter{Tổng kết}
\label{chap:ketluan}

\section{Kết quả đạt được}
% TODO: đối chiếu mục tiêu MT1--MT7 (mục~\ref{sec:muctieu}): mục tiêu nào đã hoàn thành, mức độ, kèm số liệu từ Chương~\ref{chap:danhgia}.

\section{Hạn chế của phiên bản hiện tại}

Ngoài giới hạn kỹ thuật ở mục~\ref{sec:phamvi}, một số hạng mục được chủ động đưa ra khỏi phạm vi vì không kiểm chứng được bằng số đo (theo nguyên tắc ở mục~\ref{sec:muctieudanhgia}):
\begin{itemize}
    \item \textbf{Phát hiện bất thường bằng học máy:} hiện chỉ dùng luật ngưỡng; học máy cần lưu lượng giao dịch thật đủ lớn và tập nhãn gian lận đã xác nhận, cả hai đều chưa có.
    \item \textbf{Phân quyền nội bộ trong tài khoản doanh nghiệp} (người đặt hàng, kế toán, người duyệt chi): không mang đặc thù nông sản.
    \item \textbf{Mua trả sau và hạn mức tín dụng cho khách sỉ:} kéo theo thẩm định tín dụng và thu hồi nợ, vượt phạm vi một đồ án phần mềm.
    \item \textbf{Ứng dụng di động riêng và kênh hỗ trợ qua mạng xã hội} cho nhà cung cấp ít rành công nghệ.
    \item \textbf{Luồng bán sỉ tối giản:} thương lượng nhiều vòng, giao nhiều đợt, ký quỹ và khấu trừ thuế thay chỉ dừng ở thiết kế.
\end{itemize}

\section{Hướng phát triển}

\subsection{Mở rộng về nghiệp vụ}

Mỗi hạng mục dưới đây có điều kiện kích hoạt riêng; phần mới dừng ở thiết kế được liệt kê ở Bảng~\ref{tab:chi-thiet-ke}.
\begin{itemize}
    \item \textbf{Kết nối thật} cổng thanh toán VNPay (kèm IPN) và GHN (kèm webhook) khi có tài khoản và hợp đồng; Web Push, cảnh báo chứng nhận và tồn thấp, trao đổi theo đơn, báo hàng về, thương lượng sỉ nhiều vòng; đưa dự báo giá lên màn hình thu mua khi MAPE ổn định trên dữ liệu thật.
    \item \textbf{Mô hình học máy cho giám sát bất thường} khi đã tích lũy đủ dữ liệu và nhãn từ phán quyết của nhân sự, đi từ unsupervised sang supervised.
    \item \textbf{Mở rộng vận hành:} điều phối nhiều kho sơ chế, tối ưu tuyến khi mở rộng địa bàn; đội giao hàng riêng chỉ khi mật độ đơn đủ để chi phí cố định trên mỗi đơn thấp hơn cước 3PL; hàng đợi ghi nhật ký canh tác ngoại tuyến có tự gửi lại và chống ghi trùng bằng \texttt{Idempotency-Key}; ứng dụng di động gốc cho nhân viên kho và hỗ trợ vùng khi ứng dụng web không đáp ứng nhu cầu ngoại tuyến.
    \item \textbf{Mở rộng kinh doanh:} dịch vụ kho mở cho nhà cung cấp rồi doanh nghiệp ngoài sàn; chào sỉ phần tồn dư của Farmery khi quy tắc không cạnh tranh với nhà cung cấp đã được kiểm chứng; gói hiển thị nổi bật luôn gắn nhãn ``Tài trợ''; ký số hợp đồng sỉ; cảnh báo khan hiếm cho khách sỉ.
    \item \textbf{Xuất khẩu:} bổ sung giấy tờ theo lô (mục~\ref{subsec:chungnhan}); theo phân tích khoảng cách ở mục~\ref{subsec:nhatkycanhtac-xuatkhau}, làm theo thứ tự: thay trường vùng trồng tự do bằng tham chiếu mã số vùng trồng chính thức, đặt thời hạn lưu nhật ký theo thị trường đích, mở rộng mô hình truy xuất sang cơ sở đóng gói và đơn vị kiểm dịch.
    \item \textbf{Đồng bộ TraceViet} khi được cấp quyền: đẩy dữ liệu lô đã duyệt theo chuẩn đã công bố, hiển thị mã truy xuất do Bộ cấp; khi có giao diện tra cứu mã vùng trồng và chứng nhận cho bên ngoài thì thay đối chiếu thủ công.
\end{itemize}

\subsection{Mở rộng về kiến trúc}

Modular monolith (mục~\ref{sec:kientruc}) có thể tách dần nhờ ba ràng buộc RB1--RB3. Thứ tự hợp lý: (1) trang truy xuất công khai, vì chỉ đọc, không nằm trong giao dịch và dễ bị tải đột biến khi nhiều người quét cùng lúc; (2) thông báo, vì đã chạy qua bus sự kiện (khi tách thì thay bằng hàng đợi bền như Redis); (3) tìm kiếm, chuyển sang search engine chuyên dụng khi thời gian truy vấn vượt ngưỡng NFR1.1. Ngược lại \textbf{tồn kho và đơn hàng nên giữ chung} càng lâu càng tốt vì chống bán vượt đang dựa vào transaction có row-level locking trong một cơ sở dữ liệu; tách ra phải thay bằng giữ chỗ và giao dịch bù phức tạp hơn nhiều. Với quy mô Farmery hướng tới, nhiều khả năng việc tách sẽ không bao giờ cần: microservices chỉ đáng khi có nhiều đội phát triển song song, không phải khi cần hệ thống chạy nhanh hơn.
